\section{Curve functionals: properties, general context and examples}
\label{sec:functionals}




In this section we state the main technical hypotheses that we require of our curve functionals. 

Then, we give examples and non-examples of curve
functionals satisfying the hypotheses of Theorem~\ref{thm:intersect}.
In particular, we discuss how our result fits into our previous work
in~\cite{MGT21:Smoothings}.
We then turn to further questions motivated by our results.

\subsection{Curve functionals and their properties}
\label{subsec:mainhypotheses}

We now describe in detail the hypotheses of Theorem~\ref{thm:intersect}.



In the sequel, whenever we say multi-curves, we mean oriented multi-curves, unless stated explicitly otherwise, and denote the space by $\Curves^+(S)$. In terms of geodesic currents, this means we will work with oriented geodesic currents (see Section~\ref{sec:currents} for definitions).

In this paper, the geometric intersection number between curves will play a prominent role. In general, when one considers geometric intersection number $i(C, D)$ between two multi-curves $C$ and $D$, it is assumed that $C,D$ are unoriented. If $C$ and $D$ are oriented, and $C',D'$ are different oriented multi-curves with $C'$ (resp. $D'$) equivalent to $C$ (resp. $D$)  as unoriented multi-curves, we declare $i(C',D')=i(C,D)$, in other words, define the intersection for oriented curves by forgetting the orientation (see~Definition~\ref{def:intersection_currents}).

As we shall see in the Appendix~\ref{sec:asymmetric}, there are versions of
intersection numbers for oriented curves which are asymmetric (i.e. they distinguish the
orientations of the intersections). For those, the analog of the main duality
result, Theorem~\ref{thm:intersect}, will not hold. This will also be reflected in the proof (see Proposition~\ref{prop:recoverintersection} and the related Remark~\ref{rmk:use_symmetry}).
However, in a lot of our examples the curve functionals will not distinguish orientation. Equivalently, they can be seen as curve functionals on unoriented multi-curves.
This is a natural perspective since many currents, such as most measured laminations, are unoriented.


\begin{definition}\label{def:properties}
  Let $S$ be a \emph{closed} surface, i.e., a compact topological surface without boundary,
  and
  $f \colon \Curves^+(S) \to \bbR$ be a
  function defined on the space of oriented,
  not-necessarily-simple multi-curves.
  We will also refer to $f$ as a \emph{curve functional} for short.
  (\emph{Functional} means that it takes values in scalars; it is not
  assumed to be linear, although the relevant examples to which our Theorem~\ref{thm:intersect} will apply will be additive).
  We define several
  properties that $f$ might satisfy.

  \begin{itemize}
    \item \textbf{Symmetry:} Let $C$ be an oriented multi-curve, and $C^{-1}$ the multi-curve where all the components have the opposite orientation. We have $f(C)=f(C^{-1})$.
        \item \textbf{Orientation invariance:} Let $C$ be an oriented closed multi-curve, and $C'$ the same unoriented multi-curve with a different choice of orientation. We have $f(C)=f(C')$. This condition is stronger than symmetry, but if $f$ is additive (see below), they are equivalent.
    \item \textbf{Oriented Smoothing:} Let $C$~be an oriented curve on~$S$, and
    let $x$~be an essential crossing of~$C$. Let $C'$ be the oriented
    smoothing of~$C$ at~$x$. Then $f(C) \ge f(C')$. Schematically, we
    have
    \begin{equation}\label{eq:smoothing}
      f\left(\mfig{curves-3}\right) \ge  f\left(\mfig{curves-4}\right).
    \end{equation}
    See Definition~\ref{def:essential-crossing} for ``essential
    crossing''. Loosely, it is a crossing that cannot be removed by homotopy. 
  \item \textbf{Additivity: } Let $C_1$ and~$C_2$ be two oriented curves
    on~$S$. Then
    \begin{equation}\label{eq:union-additive}
      f(C_1 \cup C_2) = f(C_1) + f(C_2).
    \end{equation}
  \end{itemize}
\end{definition}

The oriented smoothing condition will be further divided into three
types, following the convention in
Definition~\ref{def:smoothing_types}, depending on if the essential
crossing is between two connected components of the multi-curve
(\emph{oriented disconnected smoothing}); if it is a self-crossing of
the one connected component not
coming from a power of a single element (\emph{oriented connected
  smoothing}); or
if it is a self-crossing arising from a power (\emph{power
  smoothing}). Accordingly, we will
say $f$ satisfies a particular type of smoothing property if $f$ is
non-increasing under such type.

Let $[a]$ denote the conjugacy class of $a \in \pi_1(S)$, which, equivalently, can be thought of as the free homotopy class of the oriented curve represented by $a$. For $a, b\in \pi_1(S)$, let $[a][b]$ denote the multi-curve $[a] \cup [b]$. 

\begin{lemma} Oriented smoothing and oriented max-smoothing can be phrased equivalently as
follows. 
\DMGtodo{added this oriented algebraic version before the unoriented one, and made it into a lemma, so that it is clear that we are claiming an equivalence and not just stating another definition}
\begin{itemize} \item (oriented disconnected smoothing) For any pair of
  elements $a,b \in \pi_1(S)$ whose hyperbolic axes intersect, we have
  \[ f([a] \cup [b]) \geq f([ab]), \]
  where $B$ denotes the inverse of $b$; if the inequality is
  strengthened to an equality we get \emph{oriented disconnected max-smoothing}.
\item (oriented connected smoothing) for any pair of elements $a,b \in
  \pi_1(S)$ whose hyperbolic axes are parallel (i.e., do not intersect
  and are oriented coherently), we have
  \[ f([ab]) \geq  f([a] \cup [b]); \]
  similarly, if the inequality is
  strengthened to an equality we get \emph{oriented connected max-smoothing}.
\item (power smoothing) and for any element $a \in \pi_1(S)$, and any
  integers $n,m\geq 1$, we have \begin{equation} f([a^{n+m}]) \geq f([a^{ m}]) +
  f([a^{ n}]).\label{eq:power_smoothing}\end{equation}
\end{itemize}
\label{lem:oriented_smoothing_lemma_top_alg}
\end{lemma}


\begin{definition}
  For $C$ an oriented multi-curve,  $C^n$ is the oriented multi-curve with as many components as~$C$, in which each
  component of~$C$ is covered by an $n$-fold cover. That is, if
  $g \in \pi_1(S,x)$ represents $C$, $g^n$
  represents~$C^n$.
  On the other hand, $nC$ is the oriented multi-curve
  that consists of $n$~parallel
  copies of~$C$ (so with $n$ times as many components) or,
  in the context of weighted oriented multi-curves,
  with weights multiplied by~$n$),
\end{definition}

\begin{definition}\label{def:stable}
     Let $f$ be a curve
     functional and let $n>0$ be an
     integer. We say $f$ satisfies \emph{stability} if, for an arbitrary oriented multi-curve~$C$,
       \begin{equation}
         \label{eq:stable}
         f(C^n) = f(nC).
       \end{equation}
       Note that stability implies $f$ satisfies power smoothing.
\end{definition}

\begin{definition}\label{def:homogeneous} For an arbitrary oriented multi-curve~$C$, we say $f$ satisfies \emph{homogeneity}
       \begin{equation}
         \label{eq:homogeneous}
         f(nC) = nf(C).
       \end{equation}
       Note that functionals satisfying additivity (Equation~\ref{eq:union-additive}) automatically satisfy homogeneity.
\end{definition}

We emphasize that oriented connected smoothing does not imply oriented disconnected smoothing, or viceversa. 
For example, the curve functional $f$ assigning to a curve its number of connected components is a curve functional satisfying oriented disconnected smoothing but not oriented connected smoothing. Similarly, the negative of this functional satisfies oriented connected smoothing but not oriented disconnected smoothing. This example is not a stable function, but a bit more work produces examples satisfying additivity and stability while exhibiting exactly one of the smoothing properties. For example, in Lemma~\ref{lem:red_not_smoothing} we give examples of lengths coming from reducible actions of surface groups on real trees, which satisfy disconnected smoothing but not connected smoothing and viceversa. Similarly, in Example~\ref{ex:linEL} we present an additive variation of extremal length that satisfies stability, disconnected smoothing but not connected smoothing. 
\DMGtodo{this paragraph has been edited and expanded, and added a bit more detail about optimality of conditions}
One can also find examples of functions satisfying stability, and smoothing but not additivity (such as extremal length, as we will discuss in Example~\ref{ex:linEL}).
Finally, the proof of Theorem~\ref{thm:intersections_group} shows, in particular, that smoothing  implies stability. 

There exist slight generalizations of both smothing as well as additivity that have been used in our previous work~\cite{MGT21:Smoothings}.

\begin{definition}
\begin{itemize}
    \item \textbf{Oriented Quasi-smoothing:} There is a constant $R\ge 0$ with the
      following property. Let $C$~be an oriented curve on~$S$, and
    let $x$~be an essential crossing of~$C$. Let $C'$ be the oriented
    smoothing of~$C$ at~$x$. Then $f(C) \ge f(C')-R$. Schematically, we
    have
    \begin{equation}\label{eq:quasismoothing}
      f\left(\mfig{curves-3}\right) \ge f\left(\mfig{curves-4}\right) -R.
    \end{equation}
    See Definition~\ref{def:essential-crossing} for ``essential
    crossing''. If $R=0$, we recover the oriented smoothing property. 
    
    We note that the above properties are what in~\cite[Equations~(1.2),(1.3)]{MGT21:Smoothings} we called \emph{quasi-smoothing} and \emph{smoothing}. In this paper we distinguish them from their unoriented counterparts (see below). \DMGtodo{added this remark about the relation to the other paper}
  \item \textbf{Convex union:} Let $C_1$ and~$C_2$ be two oriented curves
    on~$S$. Then
    \begin{equation}\label{eq:union-convex}
      f(C_1 \cup C_2) \le f(C_1) + f(C_2).
    \end{equation}
    If the above inequalities are strengthened to equalities, we recover the additivity property (Equation~\ref{eq:union-additive}).
\end{itemize}
\end{definition}




Theorem~\ref{thm:intersect} considered curve functions on the space of unoriented multi-curves $\Curves(S)$.

\begin{definition}[smoothing and max-smoothing]
Let $C$ be an unoriented closed curve, and $C'$ any of the two smoothings, oriented or unoriented. We say $f \colon \Curves(S) \to \mathbb{R}$ satisfies \emph{smoothing} if $f(C) \geq f(C')$, Schematically, we have
     \[
      f\left(\mfig{curves-0}\right) \ge \max \left\{ f\left(\mfig{curves-1}\right), f\left(\mfig{curves-2}\right)\right\}.
   \]
If all these inequalities are strengthened to equalities, we will refer to the property as \emph{max-smoothing}.
\label{def:smoothing_and_max}
\end{definition}

Smoothing is satisfied by many natural functions, including the intersection number $i(\mu, \cdot)$ with any geodesic current $\mu$ (see~\cite[Section~4.1]{MGT21:Smoothings} cf. Appendix~\ref{sec:necessary}),
or curve lengths coming from a length metric structure on $S$~~\cite[Section~4.3]{MGT21:Smoothings}, cf. Section~\ref{subsec:domination}). We will see many more examples in Section~\ref{subsec:examples_and_non} below.

 Similarly as before, smoothing and max-smoothing can be phrased more algebraically. 
 \begin{definition}[smoothing for oriented curve functionals] \DMGtodo{note: this is NOT the same as oriented smoothing}
 Given an \emph{oriented} curve functional $f \colon \Curves^+(S) \to \mathbb{R}$, we will say it satisfies \emph{smoothing} if it satisfies the following hypotheses:
\begin{itemize} \item (disconnected smoothing) For any pair of
  elements $a,b \in \pi_1(S)$ whose hyperbolic axes intersect, we have 
  \[ f([a] [b]) \geq \max\{\,f([ab]), f([aB])\,\}; \] if the inequality is
  strengthened to an equality we obtain \emph{disconnected max-smoothing}.
\item (connected smoothing) for any pair of elements $a,b \in
  \pi_1(S)$ whose hyperbolic axes are parallel, we have
  \[ f([ab]) \geq \max\{\,f([aB]), f([a] [b])\,\}; \]
  if the inequality is
  strengthened to an equality we get \emph{connected max-smoothing}.
 \item (power smoothing) and for any element $a \in \pi_1(S)$, and any
  integers $n,m\geq 1$, we have \begin{equation} f([a^{n+m}]) \geq f([a^{ m}]) +
  f([a^{ n}]).\label{eq:power_smoothing_unoriented}\end{equation}
\end{itemize}
\label{def:smoothing_for_oriented}
\end{definition}

\DMGtodo{edited this, since I added the algebraic oriented version before}

Given a curve functional on unoriented multi-curves $f \colon \Curves(S) \to \mathbb{R}$ one can always extend it to a functional on oriented multi-curves by forgetting orientation:
\[
f^+(C) \coloneqq f(\overline{C})
\]
where $\overline{C}$ denotes the unoriented multi-curve represented by $C$.
Also, if $f \colon \Curves^+(S) \to \mathbb{R}$ is an orientation invariant oriented curve functional, then it induces a curve functional on unoriented multi-curves $\overline{f} \colon \Curves(S) \to \mathbb{R}$ given by $\bar{f}(\overline{C}) \coloneqq f(C)$--- where $\overline{C}$ denotes the unoriented curve represented by $C$.

\DMGtodo{stated the explicit equivalence}
\begin{lemma} We have the following equivalences.
\begin{itemize} \item An unoriented curve functional $f \colon \Curves(S) \to \mathbb{R}$ satisfies smoothing (resp. max-smoothing) according to Def.~\ref{def:smoothing_and_max} if and only if $f^{+}$ satisfies smoothing (resp. max-smoothing) according to Def.~\ref{def:smoothing_for_oriented}.
\item An oriented curve functional $f \colon \Curves^+(S)\to \mathbb{R}$ which is orientation invariant satisfies smoothing (resp. max-smoothing) according to Def.~\ref{def:smoothing_for_oriented}. if and only if $\overline{f}$ satisfies smoothing (resp. max-smoothing) according to Def.~\ref{def:smoothing_and_max}.
\end{itemize}
\label{lem:smoothing_lemma_top_alg}
\end{lemma}

Similarly, the above equivalences can be stated for quasi-smoothing.

We refer the reader to Section~\ref{subsec:crossings} for precise
definitions and proofs of Lemma~\ref{lem:oriented_smoothing_lemma_top_alg} and~\ref{lem:smoothing_lemma_top_alg}. See in particular
Proposition~\ref{prop:crossings-triality}.

\subsection{General context and examples}

Here we discuss how the results in this paper fit into a broader context initiated in~\cite{MGT21:Smoothings}, discuss how it relates to other recent results in the literature, and give examples of curve functionals that satisfy and do not satisfy the hypotheses of the main theorem.
 This part is independent from the rest of the paper, and is not
used in the proof in the main results.


\begin{definition}[Quasi-convex, convex, and additive-convex curve functionals]
In decreasing order of generality, we consider the following subsets
of curve functionals.
\begin{enumerate}
    \item Given $R>0$, let $\QConvex_R^+(S)$ be the set of 
      curve functionals satisfying $R$-oriented quasi-smoothing, convex union,
      stability and homogeneity. We call this the space of
      \emph{oriented $R$-quasi-convex functionals}. Define $\QConvex^+(S)
      \coloneqq \bigcup_{R \in \mathbb{R}_+} \QConvex_R^+(S)$, the space of
      \emph{oriented quasiconvex curve functionals}.
    \item Given $R>0$, we can define the subset $\QConvex_R(S)$ be the set of symmetric
      curve functionals in $\QConvex_R^+(S)$ further satisfying symmetry and $R$-quasi-smoothing (both oriented and unoriented). We call this the space of
      \emph{$R$-quasi-convex functionals} and define $\QConvex(S)
      \coloneqq \bigcup_{R \in \mathbb{R}_+} \QConvex_R(S)$, the space of
      \emph{quasiconvex curve functionals}.
    \item Let $\Convex(S) \coloneqq \QConvex_{0}(S)$ be the space of \emph{convex curve functionals}.
    (The ``Convex'' is because such functions induce convex functions on train track charts; see Section~\ref{subsubsec:convexity_charts} below.)
    \item Let $\AConvex(S)$ be the space of \emph{additive convex curve functionals}, i.e., the subset of $\Convex(S)$ of curve functionals that moreover satisfy the additive property.
    (Here the ``A'' is for additivity).
\end{enumerate}
\end{definition}

Examples below show we have strict inclusions
\[
\AConvex(S) \subsetneq \Convex(S) \subsetneq  \QConvex(S) \subsetneq \QConvex^+(S)
\]



\subsection{Examples and non-examples}
\label{subsec:examples_and_non}
In previous work, we proved the following result.

\begin{theorem}{\cite[Theorem~A]{MGT21:Smoothings}}\label{thm:convex}
   Let $f \in \QConvex^+(S)$. Then there is a unique continuous homogeneous function
  $\bar{f} \colon \Curr^+(S) \to \bbR_{\ge 0}$ that extends~$f$.
\end{theorem}

It is easy to see that the conditions of symmetry, smoothing, additivity and
stability, defining the space $\AConvex(S)$ are satisfied for any functional $i(\mu,\cdot) \colon \Curves^+(S)
\to \mathbb{R}$, for $\mu$ a geodesic current. This was proven
in \cite[Section~4]{MGT21:Smoothings}: verifying them when $\mu$ is a multi-curve and then extending to any $\mu$ using density of weighted curves, and continuity and linearity of intersection number. \DMGtodo{added this reminder about how the proof goes} We also give a different proof more in line of the types of argument in this paper in Section~\ref{sec:necessary}.
Many natural notions of length do satisfy all these properties and they have been shown to extend continuously to geodesic currents as intersection numbers. We list them here, along with the appropriate references.

\textbf{Examples}

\begin{enumerate}
   \item Intersection number with a fixed geodesic current. See \cite[Proposition~4.5]{Bonahon86:EndsHyperbolicManifolds}.
    \item Hyperbolic lengths  \cite[Proposition~14]{Bonahon88:GeodesicCurrent}.
    \item More generally, length w.r.t.\ negatively curved Riemannian metrics \cite[Th\'eor\`eme~1]{Otal90:SpectreMarqueNegative}.
    \item Even more generally, non-positively curved Riemannian metrics
\cite[Theorem~A]{CFF92:RigidityNonPosCurvedRiem}.
\item Lengths with respect to singular Euclidean metrics coming from quadratic differentials \cite[Lemma~9]{DLR10:DegenerationFlatMetrics}.
\item More generally, lengths with respect to negatively curved Riemannian  metrics with conical singularities of angle $\geq 2\pi$ \cite[Theorem~A]{HP97:RigidityNegCurvedCone}.
\item Non-positively curved Euclidean cone metrics of angle $\geq 2\pi$ \cite[Proposition~3.3]{BL17:RigidityFlat}, \cite{Con18:MarkedNonpos}.
    \item Word length with respect to simple generating sets \cite[Theorem~1.2]{Erlandsson:WordLength}.
    \item Lengths associated to Anosov representations induced by positive generalized cross-ratios \cite[Proposition~2.24]{MZ19:PositivelyRatioed}, \cite[Theorem~1.6]{BIPP24:PositiveCrossratios}.
\end{enumerate}

We emphasize that our main result Theorem~\ref{thm:intersect} subsumes all these previous results, as we will elucidate in Section~\ref{sec:special-classes} for items (1)--(8) and Section~\ref{sec:crossratios} for item (9).
In particular, in Section~
\ref{subsec:embedded_graphs} we will show that item (8) can be further generalized to stable lengths associated to embedded graphs on the surface;  \DMGtodo{it would be nice to mention an example of an embedded graph which is not a simple gen set in that section (section 8)} and in Section~\ref{subsec:ineqintersection}, we discuss how any curve length arising from a length metric structure on $S$ induces a dual geodesic current.

\textbf{Non-examples}

In order to illustrate the complexity of the problem,
we now show examples of many natural functions on curves that do \emph{not} satisfy the $\AConvex$ properties (even though many of them still extend continuously to geodesic currents).

\begin{enumerate}
    \item Some curve functionals are not symmetric, but satisfy oriented quasi-smoothing, giving examples in $\QConvex^+(S)$ (see also item (3)). For instance, the examples discussed in~\cite[Section~4.6]{MGT21:Smoothings}, stable lengths of Green metrics coming from asymmetric from generating sets (see Section~\ref{subsec:green}), or natural oriented length functions coming from Anosov representations~\cite{BridgemanCanaryLabourieSambarino18:SimpleRoots}.  
    \item Square root of extremal length, although it satisfies symmetry, smoothing, stability, convex union and homogeneity (giving an example in $\Convex(S)$), does not satisfy additive union (see \cite[Section~4.8]{MGT21:Smoothings}). Functionals of such class extend to geodesic currents by~\cite[Theorem~A]{MGT21:Smoothings}.
    It's an example of curve functional illustrating the proper containment \[\AConvex(S) \subsetneq \Convex(S).\]
    \item Although word length w.r.t.\ an arbitrary generating set of $\pi_1$ satisfies additivity (just by linear extension), it does not always satisfy smoothing nor stability (see~\cite[Example~4.10]{MGT21:Smoothings},~\cite[Theorem~1.5]{EPS20:CountingCurves},~\cite{Erlandsson:WordLength}). The stabilized version does not always satisfy smoothing, either (see the same example and also \cite[Theorem~B]{MGT21:Smoothings}). This illustrates the proper inclusion $\Convex(S) \subsetneq  \QConvex_R(S)$ for $R>0$.
    \item Another example of the above are stable lengths associated to quasi-fuchsian representations. They satisfy quasi-smoothing by \cite[Lemma~4.3]{BridgemanTaylor08:ExtensionWP} or \cite[Proposition 4.11]{MGT21:Smoothings}, and they satisfy strict smoothing if only if they are fuchsian \cite[Theorem A]{FF22:QF}, illustrating another example of the  proper inclusion $\Convex(S) \subsetneq  \QConvex_R(S)$ for $R>0$.  
    \begin{question}
    If $K$ is the quasi-fuchsian constant of the limit set of a quasi-fuchsian representation $\rho$, and $R>0$ is the minimal constant so that $\ell_{\rho} \in \QConvex_R(S)$, do we have that $R \to \infty$ as $K \to \infty$? 
    \end{question} 
\item The previous example falls in the more general framework of length functions coming from higher-rank representations of $\pi_1(S)$ into connected real semisimple algebraic group of non-compact type $G$, such as $\SL(n,\mathbb{R})$.
Martone-Zhang~\cite{MZ19:PositivelyRatioed} proved that a certain class of such representations, called \emph{positively ratioed}, induce curve functionals dual to geodesic currents, and furthermore showed they satisfy some of the inequalities in smoothing. In fact, by Theorem~\ref{thm:intersect}, they satisfy smoothing if and only if they are positively ratioed. 
        \item  Translation lengths associated to biautomatic structures. Automatic structures provide regular-language normal forms for elements of a group, yielding in efficient solutions to the word problem. Many groups of geometric origin—such as hyperbolic groups—admit automatic structures. A biautomatic structure is an automatic structure satisfying additional fellow–travelling properties ensuring good compatibility with both left and right multiplication. These stronger conditions encode deeper geometric control on the Cayley graph and imply, for instance, solvability of the conjugacy problem. We commend the reader to~\cite[Section~2.2]{HV24:Biautomaticity} for details.
    Given a biautomatic structure $\mathcal{L}$ on a surface group $\pi_1(S)$, there is a notion of translation length $\| \cdot \|_{\mathcal{L}}$ coming from its defining regular language. Hughes-Valiunas~\cite[Proposition~4.2]{HV24:Biautomaticity} have shown these translation lengths satisfy quasi-smoothing and used this to construct a new example of hierarchically hyperbolic group that is not biautomatic. These structures encompass examples of word length in the previous item, so they do not satisfy smoothing in general.
    \item Square root of self-intersection number satisfies stability,
      homogeneity and smoothing, but does not satisfy additivity or
      even convex union (see~\cite[Example~1.12]{MGT21:Smoothings}).
      This is an example of a curve functional which is not in
      $\QConvex(S)$, even though it does extend to geodesic currents
      by work of Bonahon~\cite{Bonahon86:EndsHyperbolicManifolds}.
\item Translation lengths associated to embedded graphs on a surface whose actions are not proper \cite[Proposition 11]{Bonahon91:CurrentsGroups} do not satisfy quasi-smoothing in general and, moreover, do not extend continuously to geodesic currents. 
Suppose $\Gamma$ is acting on a geodesic space $X$ by isometries, and let $\ell$ be the stable length of the action. Suppose $\ell$ is non-trivial. Let $N$ be an infinite normal subgroup of $\Gamma=\pi_1(S)$ so that $\ell(b)=0$ for all $b \in N$. Let $a \in G$, with $\ell(a)>0$ and fix $b\in N$. Then $\frac{1}{2n} a^n b a^{-n} b$ converges to $a$ in the weak$^*$ topology, but since $a^n b a^{-n} b\in N$, we have $\ell(a^n b a^{-n} b)=0$.
Hence, $\ell$ cannot extend continuously. In fact, $\ell$ does not satisfy smoothing. Indeed, we can find $b \in N$ so that $\ora{a}$ and $\ora{b}$ cross. If smoothing was satisfied, by Proposition~\ref{prop:join-split-geom}, Equation~\ref{eq:cross-cancel}, there would exist $s'>0$ so that for all $n \geq s'$, we would have
\[
\ell(ba^n ba^{-n}) \geq \ell(a^{2n}).
\]
But this inequality is violated, since on the one hand, the left hand side is $0$, and, on the other hand, the right hand side is non-zero.
\item Lorentzian length of an affine action of a Fuchsian
  group~\cite{GM00:FlatLorentzian} doesn't satisfy quasi-smoothing. 
  Indeed, it satisfies symmetry, stability and additivity~\cite[p.
  1053]{GML09:ProperAffine}, but it takes both positive and negative
  values. Proposition~\ref{prop:curve-positive2} below will show this implies it cannot satisfy quasi-smoothing. Nevertheless, it extends continuously to geodesic currents,
  as shown in~\cite[6.1]{GML09:ProperAffine}. Strictly speaking, the
  paper deals with the case of surfaces with boundary, but the same
  arguments go through for closed surfaces, by~\cite{Lab01:Diffusion}.
\end{enumerate}

For an example of a curve functional that does not extend continuously to geodesic currents (and thus does not satisfy smoothing, by~\cite[Theorem~A]{MGT21:Smoothings}), we have the following.
\begin{example}\label{ex:linEL}
Let $\sqrt{EL}$ denote the (square root of) extremal length.
In \cite[Examples]{MGT21:Smoothings} the authors proved that $\sqrt{EL}$ satisfies stability, homogeneity, smoothing, and convex union. By  \cite[Theorem~A]{MGT21:Smoothings}, $\sqrt{EL}$ extends continuously to a continuous on geodesic currents.
$\sqrt{EL}$ does not satisfy additivity, so it is not dual to a geodesic current. However, $\sqrt{EL}$, as a function on connected curves can be linearly extended to a function $G$ on (weighted) multi-curves, 
(which is thus not equal to $\sqrt{EL}$ on multi-curves). $G$ cannot extend continuously to currents, since we know $G=\sqrt{EL}$ on (weighted) connected curves, which are dense in the space of geodesic currents.
Thus, the purported continuous extension $\overline{G}$ would have to satisfy $\overline{G}=\sqrt{EL}$ on all geodesic currents, hence on all multi-curves, a contradiction.
It follows from this that $G$ cannot satisfy smoothing, since it trivially satisfies convex union (in fact, additivity) and stability.  One can also see that, specifically, it does not satisfy connected smoothing.
Indeed, let $F=\sqrt{EL}$. For a crossing pair $(a,b)$, we have
\[
F(a + b) \geq F(ab).
\] 
Then, by disconnected smoothing of $F$, we have
\[
G(a+b)=G(a) + G(b)=F(a) + F(b) \geq F(a+b) \geq F(ab)=G(ab)
\]
where we used the fact that $F$ and $G$ take the same values on closed curves; hence,
$G$ satisfies disconnected smoothing.
This leads to the conclusion that $G$ cannot satisfy connected smoothing, i.e. there must be some pair of $(a,b)$ parallel so that
$G(ab) < G(a) + G(b)$. This is indeed the case. Maxime Fortier Bourque pointed the first author to~\cite{HZ20:SwissCross} for a strategy to compute the extremal length of a figure-eight curve in a pair of pants.  Consider the extremal length of a figure-eight curve in a genus two surface $S$. Consider the pair of pants $P$ bounding $ab$, and let us think of it as obtained by gluing adjacent sides in a Swiss cross. The extremal metric has been computed numerically in~\cite[Page~35]{HZ20:SwissCross} as a function of the width $w$ of the Swiss cross. If the width is chosen to be a fourth of the total size of the Swiss cross, and the normalization $w=1$ is chosen, the extremal length of the figure-eight can be computed to be 2.4251
 (see Table in \cite[Page~36]{HZ20:SwissCross}, row $t=4$), hence its square root is $G_P(ab) \simeq 1.557273$. On the other hand, the length of each one of the resulting smoothings is $3$, while the area of the surface with the extremal metric associated to $ab$ is computed as $A=6.5977$. 
 
 Hence, $G_{S}(a) + G_S(b) =2G_S(a)=\sqrt{2}G_P(a)$ is lower bounded by $\simeq 1.6517335$, showing the inequality
 \[
 G_S(ab) < G_P(ab) < G_S(a) + G_S(b) = 2G_S(a)=\sqrt{2} G_P(a).
 \]
\end{example}


\textbf{Non-negativity}

Observe that we do not require in the hypotheses that the curve
functional $f \in \AConvex(S)$ is non-negative: in fact it follows
from the other conditions.
We finish this subsection by relating the hypotheses on $f$ to non-negativity. In particular, we prove that smoothing and convex union on $f$, imply non-negativity.

\begin{proposition}
  Let $f$ be a curve functional satisfying convex union and smoothing. Then $f([a]) \geq 0$ for any $a \in \pi_1(S)$.
  \label{prop:curve-positive}
\end{proposition}
\begin{proof}
Let $x \in \pi_1(S)$ be so that $\ora{a}$ crosses $\ora{x}$ to the
right, so
$[a][x] \reducesto [ax]$.
By Lemma~\ref{lem:rhs_prod}, we see that
$(\ora{ax},\ora{a})$ L-cross, so
(by unoriented smoothing) $[ax][a]
\reducesto [x]$. We thus have
\begin{align*}
    f([a]) + f([x]) &\geq f([a] [x]) \geq f([ax])\\
     f([ax]) + f([a]) &\geq f([ax] [a]) \geq f([x]).
\end{align*}
Adding the equations, we get $f([a]) \geq 0$.
\end{proof}

Although we will not use it in this paper, we note that a similar
argument with different assumptions also yields positivity.

\begin{proposition}
  Let $f$ be a curve functional satisfying convex union, quasi-smoothing,
  stability, and homogeneity. Then $f([a]) \geq 0$ for any $a \in
  \pi_1(S)$.
  \label{prop:curve-positive2}
\end{proposition}

\begin{proof}
Let $x \in \pi_1(S)$ be so that $\ora{a}$ crosses $\ora{x}$ to the
right. As in the proof of Proposition~\ref{prop:curve-positive}, we have
\begin{align*}
[a^n][x] &\reducesto [a^nx]\\
[a^nx][a^n]&\reducesto [x].
\end{align*}
Since $f$ satisfies quasi-smoothing (oriented and unoriented) as well as convex union, we get
\begin{align*}
    f([a^n]) + f([x]) &\geq f([a^n] [x]) \geq f([a^nx]) - R\\
     f([a^nx]) + f([a^n]) &\geq f([ax] [a^n]) \geq f([x]) - R.
\end{align*}
Adding the equations, we get $f([a^n]) \geq -2R$. By stability and
homogeneity, this yields $n f([a]) \geq -2R$, and taking $n$ large
enough we see that $f([a]) \geq 0$.
\end{proof}

\begin{remark}
The use of both the oriented and the unoriented smoothing hypothesis is essential in both of the above results. In fact, there exist natural examples of functions that satisfy the oriented quasi-smoothing condition (but not the unoriented one), as well as all the other necessary hypotheses mentioned above, yet still take negative values. For instance, curve functionals arising from homogeneous quasi-morphisms (see Section~\ref{sub:quasimorphism}).
\end{remark}


\section{Further questions}
\label{sec:questions}

We finish by discussing some natural follow-up questions and conjectures.

\subsection{Open surfaces}
A natural question is whether our results extend to open surfaces, i.e.\ compact surfaces with boundary or finite-type surfaces with punctures.

For such surfaces, geodesic currents admit at least two natural definitions, depending on how one treats the ends. 
If the ends are realized as cusps, one obtains a space $\Curr_{\mathrm{cusp}}(S)$, defined analogously to Definition~\ref{def:currents} as the space of $\pi_1(S)$-invariant 
Radon measures on the space of geodesics in the universal cover (see, e.g.,~\cite[Definition~2.2]{Sas22:Currents}). In this setting, the space of geodesics includes geodesic arcs connecting cusps. 
If instead one restricts to geodesics projecting to the convex core,
one obtains a space denoted $\Curr_{\mathrm{open}}(S)$ (\cite[Section~2.6]{DLR10:DegenerationFlatMetrics}). See
Figure~\ref{fig:open_surfaces}.

\begin{figure}
\centering
\begin{subfigure}{0.48\textwidth}
\centering
\resizebox{60mm}{!}{
\svgfigbig{cusp}}
\caption{Cusped surface with an arc $A$ and a closed geodesic $C_n$
  wrapping around the arc. We have $C_n \to 2A$.}
\end{subfigure}
\hfill
\begin{subfigure}{0.48\textwidth}
\centering
\resizebox{60mm}{!}{
\svgfigbig{open}}
\caption{Open surface with two boundary parallel curves $B_1$ and $B_2$. The sequence of curves $C_n$ also sits in this surface, and we have $\frac{1}{n}C_n \to B_1 + B_2$.}
\end{subfigure}
\caption{Examples of cusped and open surfaces, and currents on them.}
\label{fig:open_surfaces}
\end{figure}

In $\Curr_{\mathrm{cusp}}(S)$, the intersection pairing fails to be
continuous; see~\cite[Remark~2.26]{MGT21:Smoothings} and
also~\cite{Sas22:Currents, T22:ThurstonCompactification}. It is thus
hard to get started working entirely in this setting.


The case of $\Curr_{\mathrm{open}}(S)$ is different. Even though intersection number is continuous in this case, it is not true that our defining properties of $\AConvex(S)$ guarantee a dual geodesic current.
The key issue is that the intersection number is degenerate: any
geodesic current $\mu$ supported on boundary curves will have zero
intersection number with all geodesic currents $\nu \in \Curr(S)$.

To better understand the relation between the cusped and open
situation, consider Figure~\ref{fig:open_surfaces} in a surface with
two ends for simplicity.
In the cusped surface, let $A$ be the simple arc going from cusp to cusp, which is in $\Curr_{\mathrm{cusp}}(S)$.
Let $C_n$ be a sequence of closed geodesics obtained by following $A$
from near one cusp to the other, wrapping around the second cusp $n$ times,
going back along $A$, and wrapping around the first cusp $n$ times
before closing up. In the weak$^*$-topology of
$\Curr_{\mathrm{cusp}}(S)$, we have $\lim_{n\to\infty}C_n = 2A$. (See, for example,
Sasaki~\cite[Section~6]{Sas22:Currents}.)

In the $\Curr_{\mathrm{open}}$ case, let $B_1$ and $B_2$ be the
boundary parallel curves homotopic to each respective boundary
component. The sequence of curves $C_n$ can also be considered in
$\Curr_{\mathrm{open}}(S)$. Without scaling, $C_n$ does not converge
to any current, but it does projectively. In fact, $\lim_{n \to
  \infty} \frac{1}{n} C_n = B_1 + B_2$.

Consider the geometric intersection number between arcs and curves in the cusped surface and curves in the open surface, i.e., $i \colon \Curves_{\mathrm{cusp}}(S) \times \Curves_{\mathrm{open}}(S)$.
We claim
\[
i(C_n, \tfrac{1}{n} C_n) \to i(2A, B_1 + B_2).
\]
Indeed, $i(C_n, C_n) = 4n$, and $i(2A, B_1+B_2)=4$, from which the result follows.
Following this example, we conjecture there is a well-defined and
continuous geometric intersection number between
$\Curr_{\mathrm{open}}(S)$ and $\Curr_{\mathrm{cusp}}(S)$.

\begin{conjecture}
There exist a continuous function 
\[
I \colon \Curr_{\mathrm{cusp}}(S) \times \Curr_{\mathrm{open}}(S) \to \mathbb{R}
\]
extending the natural geometric intersection pairing between $\Curves_{\mathrm{cusp}}(S)$ and $\Curves_{\mathrm{open}}(S)$.
\end{conjecture}

It is known that if one restricts the right-hand-side
factor to currents supported on a preimage of a compact subset on the
surface, the intersection number is continuous
(See \cite[Theorem~2.4]{BIPP21:Systoles}
and~\cite[Section~6]{Sas22:Currents}).

We can use this intersection number to propose a duality. Say that a
function $f \colon \Curves_{\mathrm{cusp}}(S) \to \mathbb{R}$
satisfies \emph{smoothing} if it satisfies smoothing between any
closed multi-curves, as well as smoothing between arcs and curves and
between arcs, in the sense that after resolving any essential crossing
(any intersection that cannot be homotoped away), the function does
not increase.

Similarly, say that a function
$f \colon \Curves_{\mathrm{open}}(S) \to \mathbb{R}$ satisfies
smoothing if it does not increase after resolving any essential
crossing.

For examples of such functionals, consider an embedded graph $\mathcal{G} \subset S'$.
Equipping $\mathcal{G}$ with a metric, we get a notion of length of
$\ell_\Sigma$ curves in $\Curves_{\mathrm{open}}(S)$ with respect to
this embedded metric graph. (See also
Section~\ref{subsec:embedded_graphs}.) When $\mathcal{G}$ is a rose-graph
(embedded and
with one vertex), Erlandsson~\cite{Erlandsson:WordLength} has
constructed examples of geodesic currents $\mu_{\mathcal{G}}$ in
$\Curr_{\mathrm{open}}(S)$ (in fact, half-integer weighted arcs and
multi-curves) so that
\[
\ell_{\Sigma}(C)=I(\mu_{\mathcal{G}}, C).
\]
This motivates the following question.

\begin{question}\label{quest:open-to-cusp}
Given a function $f \colon \Curves_{\mathrm{open}}(S) \to \mathbb{R}$
which is additive and satisfies smoothing, are there additional
assumption guaranteeing the existence of a unique geodesic current $\mu \in f \colon \Curr_{\mathrm{cusp}}(S)$ so that
\[
f(C)=i(C, \mu)
\]
for every $C \in \Curves_{\mathrm{open}}$?
\end{question}

Unfortunately, additivity and smoothing alone are not enough. For
example, take the function
$f \colon \Curves_{\mathrm{open}}(S) \to \mathbb{R}$ assigning
hyperbolic length to non-boundary parallel curves and $0$ to
boundary-parallel curves. This functional still satisfies smoothing
(since no other curve can cross a boundary-parallel curve essentially)
but does not have a dual geodesic current in
$\Curr_{\mathrm{cusp}}(S)$, since by continuity of $I$ it would have
to assign positive length to limits of curves wrapping around the
boundary. See~\cite[Remark~2.26]{MGT21:Smoothings} for details.

In the other direction, we have a parallel conjecture.

\begin{conjecture}\label{conj:cusp-to-open}
Given a function $f \colon \Curves_{\mathrm{cusp}}(S) \to \mathbb{R}$ which is additive and satisfies smoothing, there exists a unique geodesic current $\mu \in f \colon \Curr_{\mathrm{open}}(S)$ so that
\[
f(C)=i(C, \mu)
\]
for every $C \in \Curves_{\mathrm{cusp}}$(S).
\end{conjecture}

We have two motivating examples of functionals for
Conjecture~\ref{conj:cusp-to-open}.

The first one is related to the example on graphs above. Take again an
embedded metric graph $\mathcal{G} \subset S$, and define the
intersection $i(\mathcal{G},C)$ to be the minimal transversal
intersection between homotopy classes of embedded representatives of
$\mathcal{G}$ into $S$ and a homotopic representative of $C$, for
every closed curve in $\Curves_{\mathrm{cusp}}(S)$. This defines a
function $f \colon \Curves_{\mathrm{cusp}}(S) \to \mathbb{R}$. Then
our conjecture would assert the existence of a geodesic current
$\mu \in \Curr_{\mathrm{open}}(S)$. Presumably, as in Erlandsson's
result, $\mu$ in this case consists of half-integral combination of
curves.

As a second example, consider a flat metric on the completion of $S'$ so that the missing points become conical singular points. Now the length of an arc between cusps is the length of the arc from cone point to cone point, which is a finite number.
This gives a well-defined notion of length $\ell \colon \Curves_{\mathrm{cusp}}(S) \to \mathbb{R}$.
On the other hand,  $\Curr_{\mathrm{open}}(S)$ are geodesic currents
whose support is contained in the closure of the subspace of geodesics
that do not go through cone points. The natural Liouville metric for
singular flat metrics is an example of such a
current~\cite[Section~3.2]{BL17:RigidityFlat},
and seems to be dual
to~$\ell$.


\subsection{Generalizations to hyperbolic groups}
In recent work of Cantrell and Reyes and, independently, Kapovich and the first author, the following property has played a key role in extending stable lengths coming from group actions of hyperbolic groups $\Gamma$ acting on uniformly quasi-geodesic hyperbolic spaces $X$~\cite{KMG:BBP,CR25:Manhattan, CMGR26:GreenMetrics}.
Roughly, a Gromov hyperbolic space is a metric space whose triangles are thin. A group $\Gamma$ is said to be Gromov hyperbolic if its Cayley graph is.

A metric space $X$ is uniformly quasi-geodesic if there exists
constants $\lambda >0$ and $\epsilon \geq 0$, so that
every two points $x, y \in X$ are connected by a
$(\lambda,\epsilon)$-quasi-geodesic. We say that an action
$\Gamma \acts X$ of a non-elementary hyperbolic group~$\Gamma$ on a uniformly
quasi-geodesic hyperbolic metric space~$(X,d)$ has the \emph{bounded
  backtracking property} (see~\cite[Section~6.1]{CR25:Manhattan}) if given an orbit map
$\iota \colon \Cay(\Gamma) \to X$, there exists a constant $C>0$ so that
for every two points $x,y \in \Cay(\Gamma)$, and geodesic $[x,y]$,
there exists some quasi-geodesic $[\iota(x),\iota(y)]$ in $X$ so that
$\iota([x,y])$ and $[\iota(x),\iota(y)])$ are $C$-Hausdorff close in $d$.

There are examples of stable lengths satisfying bounded backtracking
which do not satisfy quasi-smoothing
(see~\cite[Section~9]{CMGR26:GreenMetrics}). On the other hand, one
can ask if all functions on $\pi_1(S)$ satisfying quasi-smoothing
arise as the stable length of some action of $\pi_1(S)$ satisfying
bounded backtracking. 

\begin{question}
Let $\Gamma$ be a surface group, and let $f \colon \Gamma \to \mathbb{R}$ be a symmetric and stable functional satisfying quasi-smoothing (and hence, extending continuously to geodesic currents).
Does there exists a uniformly quasi-geodesic Gromov hyperbolic space $X$ on which $\Gamma$ acts with the bounded backtracking property and so that $\ell_X = f$, where $\ell_X$ denotes the stable length of the action $\Gamma \acts X$?
\end{question}

Motivated by this, it would be interesting to generalize the quasi-smoothing and smoothing condition to the general setting of hyperbolic groups.

\begin{question}
What is a generalization of the properties of smoothing and
quasi-smoothing for functionals on non-elementary hyperbolic groups
that makes the answer to the
above question work in that setting?
\end{question}

Related to Theorems~\ref{thm:intersections_group}
and~\ref{mainthm:skora}, it would be interesting to give a list of
axioms, in the spirit of oriented smoothing, characterizing
small actions of a non-elementary hyperbolic group $\Gamma$.

\begin{question}
\label{qu:generalize_intersections}
Let $\Gamma$ be a non-elementary Gromov hyperbolic group action on an
$\mathbb{R}$-tree $T$. Give axioms  characterizing when the action is small, i.e., has
virtually cyclic edge stabilizers.
\end{question}

\subsection{Quasi-morphisms} 
\label{sub:quasimorphism} 

The hypotheses of \emph{oriented quasi-smoothing} are reminiscent of, but weaker than, the conditions for a map 
$f \colon \pi_1(S) \to \mathbb{R}$ to be a \emph{quasi-morphism}, that is, there exists $D > 0$ such that 
\[
|f(ab) - f(a) - f(b)| < D \quad \text{for all } a,b \in \pi_1(S).
\]
A quasi-morphism is called \emph{homogeneous} if $f(a^n) = n f(a)$ for every $a \in \pi_1(S)$ and $n \ge 1$.

Quasi-morphisms have played a prominent role in the study of stable commutator length and bounded cohomology of groups~\cite{Cal09:scl,Pic97:BoundedCohomology_currents}.

It would be interesting to understand whether one can obtain homogeneous quasi-morphisms from curve functionals satisfying oriented quasi-smoothing and stability.

\begin{question}
Let $f \colon \pi_1(S) \to \mathbb{R}$ satisfy oriented quasi-smoothing and stability.  
Are there additional conditions on $f$ ensuring that the function
\[
g(a) \coloneqq f(a) - f(a^{-1})
\]
is a homogeneous quasi-morphism?
\end{question} 

To get a non-trivial function~$g$, the functional~$f$ must be
asymmetric. See Appendix~\ref{sec:asymmetric} for examples.

% \subsection{Functionals valued on ordered abelian groups}

% We expect most of the results in this paper to extend, with only minor modifications, to the setting of $\Lambda$-valued functionals and $\Lambda$-valued geodesic currents, where $(\Lambda, <)$ denotes an ordered abelian group---that is, an abelian group equipped with a total order satisfying $a < b$ and $c \le d$ imply $a + c < b + d$.
% This generalization is natural in the context of the theory of $\Lambda$-trees~\cite{Sko90:GeometricAction} and has already proved fruitful in the study of compactifications of character varieties of surface-group representations into higher-rank Lie groups~\cite{BIPP24:PositiveCrossratios}.
% We plan to explore these extensions and their applications in future work.

\subsection{Green metrics}
\label{subsec:green}
Let $\Gamma=\pi_1(\Sigma) \subset \Isom^+(\mathbb{H})=\PSL_2(\mathbb{R})$ be a countable, discrete cocompact subgroup, for some hyperbolic metric $\Sigma$ on $S$.
Let $\mu$ be a finitely supported probability measure on $\Gamma$ whose support contains a symmetric generating set of $\Gamma$.
Let $w_n \coloneqq g_1 \cdots g_n$ be the random walk with respect to $\mu$, where each $(g_i)$ is identically and independently distributed with distribution $\mu$. Let $(\Gamma^{\mathbb{N}}, \mu^{\mathbb{N}})$ be the step space, and the map $\pi \colon \Gamma^{\mathbb{N}} \to \Gamma^{\mathbb{N}}$ as $\pi((g_n)_{n \in \mathbb{N}}) \coloneqq (w_n)_{n \in \mathbb{N}}$. The image of $\pi$ is equipped with the measure $(\mathbb{P}_{\mu})_*(\mu^{\mathbb{N}})$.
Define the metric \[ d_{\mu}(x,y) \coloneqq - \log \mathbb{P}_{\mu}(\exists n : w_n x =y).\]
In words, we are considering the probability that a random walk departing from $x$ will reach $y$ after some number of steps.
This metric has an associated stable length function $\ell_{\mu}(g) \coloneqq \lim_n d_{\mu}(x, g^n x)/n$, for any hyperbolic element $g \in \Gamma$, which can then be thought as a curve functional $\ell_{\mu} \colon \Curves^+(S) \to \mathbb{R}$. The following question emerged from conversations of the first author with Eduardo Reyes.

\begin{question}
Does $\ell_{\mu}$ satisfy the hypotheses of
Theorem~\ref{thm:intersect} for some admissible probability measure $\mu$ (and thus come from a geodesic current)?
\label{q:green-smoothing}
\end{question}

\begin{question}
Does $\ell_\mu$ further satisfy the
hyperbolic metric hypotheses in Theorem~\ref{thm:hyp-lam-char}(H) for some admissible
probability measure~$\mu$ (and thus come
from a hyperbolic metric)?
\label{q:Green}
\end{question}

If answers to either of these questions is negative for every $\mu$,
they would imply a negative answer to the following well-known question.

\begin{question}
Given $\Sigma$, is there a finitely-supported probability measure
$\mu$ on $\Gamma=\pi_1(\Sigma)$ so that the induced translation length
$\ell_{\mu}$ is dual to the Liouville current $\mathcal{L}_{\Sigma}$?
\label{q:singularity}
\end{question}

By work of Cantrell-Tanaka~\cite[Theorem~1.2]{CT25} and
Gouëzel-Mathéus-Maucourant \cite[Theorem~1.2]{GMM18:Entropy}, a
negative answer to this question would settle the so-called
\emph{singularity conjecture} for $\SL_2(\mathbb{R})$
(see~\cite[Conjecture, page~~259]{KLP11:Singularity}). In short, by
these results Question~\ref{q:singularity} is equivalent to asking if
there is a probability measure $\mu$ on $\Gamma$ so that the induced
`hitting measure' $\nu_{\mu}$ on $\partial \Gamma$ is absolutely
continuous with respect to the Lebesgue measure (coming from the
hyperbolic metric $\Sigma$) on $\partial \Gamma$.
Kosenko-Tiozzo~\cite{KT22:Inequality} have studied the metric $d_{\mu}$ in some cases of Fuchsian groups $\Gamma$ with centrally symmetric fundamental domains.
Even a positive answer to Question~\ref{q:Green} would be interesting, since then
Theorem~\ref{thm:intersect} would give a dual geodesic current for the
Green metric, and then one could study what types of geodesic currents
can arise.

\subsection{Smoothing vs multi-smoothing}
\label{subsec:multismoothing}
The hypotheses of smoothing in Equation~\eqref{eq:smoothing} of Theorem~\ref{thm:intersect} apply to one essential crossing at a time. There is a stronger condition than smoothing, where multiple smoothings at essential crossings are performed at once, and the function is required to be non-increasing under such operation.
We will refer to this condition as \emph{multi-smoothing}.
It is clear multi-smoothing implies smoothing.
It is easily checked that intersection numbers satisfy the multi-smoothing property (see, for example, \cite{AH21:EffectiveCounting3}).
What is not so clear is that functionals of $\AConvex(S)$ (that is, satisfying stability, homogeneity, smoothing and additivity) must satisfy multi-smoothing. This is, of course, a consequence Theorem~\ref{thm:intersect}.
Other curve functionals that are not in $\AConvex(S)$ also satisfy
multi-smoothing. For example, Arana-Herrera also shows
that extremal length satisfies
multi-smoothing \cite{AH21:EffectiveCounting3}.
This prompts the following question.

\begin{question}\label{quest:multi-smoothing}
If a curve functional satisfies homogeneity, stability, convex union,
and smoothing, does it necessarily satisfy multi-smoothing?
\end{question}

In this context, Arenas and Neumann-Coto
recently studied the related concept of \emph{taut smoothing}  \cite{ANC24:TautSmoothings}.

\subsection{Domination}
\label{subsec:domination}
In another direction from Section~\ref{subsec:multismoothing},
we say that a curve $C$ \emph{dominates} another curve $C'$ if $i(C,D) \geq
i(C',D)$ for every curve $D$. We say a functional $f$ satisfies
\emph{domination} if for
every pair of curves $(C,C')$ with $C$ dominating $C'$, we have
$f(C) \geq f(C')$.
It is easy to see that if $C'$ is obtained from $C$ by
multi-smoothing, then $C$ dominates $C'$. There are however other
pairs $(C,C')$ with $C$ dominating $C'$ but $C'$ not obtainable from
$C$ via multi-smoothing
(see~\cite[Example]{NC01:ShortestGeodesics}).
It is clear, by definition, that intersection numbers satisfy domination.
All in all, we have the following implications on a functional~$f$:
\[
  \text{Domination} \Rightarrow \text{Multi-smoothing} \Rightarrow \text{Smoothing}.
\]

From Theorem~\ref{thm:intersect} we obtain the following weaker
theorem which is nevertheless conceptually interesting.

\begin{theorem}
A curve functional satisfies domination and additivity if and only if it is dual to a geodesic current.
\end{theorem}

This can be interpreted as saying that linear maps that preserve the poset structure of the space of multi-curves correspond to intersection numbers.

For similar reasons as for the smoothing (or multi-smoothing)
condition, the additivity hypothesis is necessary, since
the square
root extremal length can be easily seen to satisfy domination,
homogeneity and stability, but not additivity, and it is not dual to
an intersection number.
This suggests the following strengthening of
Question~\ref{quest:multi-smoothing}.

\begin{question}
If a curve functional satisfies homogeneity, stability,
convex union, and smoothing, does it necessarily satisfy domination?
\end{question}

\subsection{Convexity on train-track charts and local compactness}
\label{sec:convexity_charts}

One of the reasons multi-smoothing is a useful property is that,
coupled with convex union, homogeneity, and stability, it
yields convex functions when restricted to train-track coordinates.
Precisely, let $f$ be a curve functional satisfying the aforementioned
properties, including multi-smoothing. Let $\tau$ be a maximal train-track (i.e., a
train-track all of whose complementary components are trigons), let
$V(\tau)$ be the set of admissible weights on~$\tau$, $\ML(\tau)$ be the subspace of
measured laminations carried by $\tau$, and
$\phi_{\tau} \colon V(\tau) \to \ML(\tau)$ be the induced
chart on measured laminations.
Since $f$ extends continuously to geodesic currents \cite[Theorem~A]{MGT21:Smoothings}, we can consider the composition $f_{\tau} \coloneqq f \circ \phi_{\tau} \colon V(\tau) \to \mathbb{R}$. Under these assumptions, the same proof as in~~\cite[Appendix~A]{Mir04:SimpleGeodesics}, \cite[Theorem~2]{Thu16:Research}, \cite[Theorem~4.7]{AH21:EffectiveCounting3}, shows $f_{\tau}$ is convex and homogeneous. 

These curve functionals form a subset of the functionals that extend
continuously to currents, by~\cite[Theorem~A]{MGT21:Smoothings}.
Moreover, the map $\Psi \colon \Curr(S) \to \mathbb{R}^{\Curves(S)}$
defined by $\Psi(\mu) = i(\mu, \cdot)$ is a homeomorphism onto its
image when $\mathbb{R}^{\Curves(S)}$ is equipped with the product
topology,
by~\cite[Theorem~11]{DLR10:DegenerationFlatMetrics}.
Thus, a restatement of Theorem~\ref{thm:intersect} is as follows.

\begin{theorem}
   The image of~$\Psi$ is homeomorphically $\AConvex(S) \subset \bbR^{\Curves(S)}$.
\end{theorem}

It is a consequence of this last result and the local compactness of
the space of geodesic currents (recalled as
Proposition~\ref{prop:local_compactness} below) that the space
$\AConvex(S)$ with the product topology is locally compact, which is
not obvious \textit{a priori} by looking directly at the functionals.

The sublevel set $\Lambda_f \coloneqq \{ \vec{x} \in V(\tau)
: f_{\tau}(\vec{x}) \leq 1 \}$ is a compact convex subset of
$\mathbb{R}^N$, for some $N>0$ depending on the combinatorics of~$\tau$.
For any sequence of geodesic currents $\mu_i$ limiting to~$\mu$,
let $f_i(\cdot) = i(\mu_i, \cdot)$ for $f(\cdot) = i(\mu, \cdot)$.
Then $\Lambda_{f_i} \to \Lambda_f$ in the Hausdorff topology. With
respect to this topology,
the subset of compact convex subsets of $\mathbb{R}^N$ is locally
compact \cite[Chapter~6]{Lay92:Convex}. We see this phenomenon as a
shadow of the local compactness of $\AConvex(S)$, which suggests that
the same local compactness for $\Convex(S)$ should hold.

\begin{conjecture}\label{conj:convex-compact}
The set $\Convex(S)$ with the product topology is locally compact.
\end{conjecture}

Some examples of interest for this conjecture would be extremal
lengths associated to conformal structures or to elastic graphs
\cite[Sections 4.8 and 4.9]{MGT21:Smoothings}.






\subsection{Analogy to Poincar\'e duality}
\label{subsec:poincare_duality}
Theorem~\ref{thm:intersect} gives a new definition of geodesic currents in terms of functionals on curves, by giving a bijection between geodesic currents and a subset of the space of curve functionals.
As mentioned in the introduction, one can view this theorem in analogy to Poincaré duality for closed surfaces.  Given a 1-cohomology class $\alpha^*$ of an $n$-dimensional manifold (represented by homomorphisms from $\pi_1(S)$ to $\mathbb{R}$), Poincaré duality associates an $(n-1)$-dimensional homology class $\alpha$ (represented by real linear combinations of $(n-1)$-cycles). Since $n=2$, these cycles can be expressed as linear combinations of simple closed curves.  

In our setting, the role of a 1-cohomology class is played by a curve
functional that satisfies smoothing, additivity, and
stability. Although it is not a homomorphism from
$\pi_1(S)$ to anything, it preserves a certain combinatorial structure
(see Section~\ref{subsec:domination}). Likewise, the role of the homology
class is taken by a geodesic current.

Our theorem can be thought of in analogy to $\alpha \cap \beta =
  \alpha^*(\beta)$, where $\cap$ denotes the cap product in homology (algebraic intersection), and thus relating cap product to the homology-cohomology pairing.

  This aligns with our theorem, which establishes a duality via the \emph{geometric} intersection number $i(\mu_f, C) = f(C)$, suggesting that the geometric intersection number plays the role of the cap product. 

  
In Poincaré duality, one also has the standard relation $\alpha^* \cup \beta^* = \alpha \cap \beta$, where $\cup$ denotes the cup product in cohomology.  
A natural question that arises then is the following.

\begin{question} What is the appropriate analog of the cup product?
  More precisely, can one define the intersection $i(\cdot,\cdot)$ of
  two geodesic
  currents directly via their dual curve functionals, and prove its
  continuity?
\end{question}

\subsection{Asymmetric analogues}
\label{sec:asymmetric-analogues}

In this paper, we characterized additive, unoriented curve functionals that
satisfy smoothing as exactly those coming from geodesic currents.
We could instead look at functionals on \emph{oriented} curves.

\begin{definition}
  Let $\AConvex^+(S)$ be the space of functionals on oriented curves
  $\Curves^+(S)$ that satisfy additivity, oriented smoothing
  (connected and disconnected), and stability, endowed with the
  product topology. Similarly, $\AConvex^+_{\ge 0}(S)$ is the subspace
  of positive functionals.
\end{definition}

Appendix~\ref{sec:asymmetric} can be summarized as saying that
$\AConvex^+(S)$ is not very related to the space $\Curr^+(S)$ of
oriented geodesic currents: the natural map $\mu \mapsto
\vec\imath(\mu,\cdot)$ is neither injective nor surjective, in what
appear to be large ways.

Nevertheless the space $\AConvex^+_{\ge 0}(S)$ may be of independent
interest. There are many examples of such functionals. For instance,
you can take an asymmetric metric on $S$ and obtain a functional
on curves as in Section~\ref{subsec:ineqintersection}. As in the
symmetric case, if there are rectifiable curves in each homotopy class
the result is automatically positive and satisfies oriented smoothing.

For a more concrete example, you can take an asymmetric generating set
for $\pi_1(S)$ as a monoid whose (directed) graph is simple in~$S$, as
in the example in Appendix~\ref{sec:asymmetric}.

\begin{question}\label{quest:AConvex+-compact}
  Is the space $\AConvex^+_{\ge 0}(S)$ a reasonable space? For
  instance, is it locally compact?
\end{question}

\begin{question}
  Does every element of $\AConvex^+_{\ge 0}(S)$ come from an
  asymmetric pseudo-metric in the style of
  Corollary~\ref{maincor:length_dual}?
\end{question}

In this oriented context, the restriction that the functional is
positive is non-trivial (unlike for unoriented functionals): for any
functional $f$ and element $\omega \in H^1(S; \bbZ)$, for a sufficiently
large constant $\lambda>0$ the combination $f + \lambda\omega$ will take negative values.

We could also relax the condition of smoothing to quasi-smoothing, and
ask similar questions. If we drop positivity and also consider
quasi-smoothings, then the space of functionals includes the much
larger space of quasi-morphisms (Section~\ref{sub:quasimorphism}), which an infinite dimensional Banach space and hence not locally compact (see, e.g.,~\cite{CRH26:Quasimorphisms}).

Recall that $\QConvex_K^+(S) \subset \QConvex^+(S)$ denotes the subspace of curve functionals on oriented multi-curves satisfying oriented quasi-smoothing with constant~$K$, homogeneity, stability and convex union.
$\QConvex_K^+(S)$ contains, for example, the space of homogeneous $K$-quasi-morphisms. The space of $K$-quasi-morphisms is easily seen to be locally compact by an inequality
\[
f(C) \leq K' \ell_F(C)
\]
where $\ell_F$ is the word-length induced by any finite generating set, and $K'$ is a constant depending only on a finite subset of curves.
The hypotheses of oriented $K$-quasi-smoothing are much weaker than those of $K$-quasi-morphisms.

\begin{question}\label{quest:QConvex+-compact}
Is $\QConvex_K^+(S)$ locally compact?
\end{question}

More generally, one could ask about local compactness for various weakenings of the the hypotheses defining the subset $\AConvex(S)$; Conjecture~\ref{conj:convex-compact} and Questions~\ref{quest:AConvex+-compact} and \ref{quest:QConvex+-compact} are some examples.

%%% Local Variables:
%%% mode: latex
%%% TeX-master: "Intersections"
%%% End:



%%% Local Variables:
%%% mode: latex
%%% TeX-master: "Intersections"
%%% End:
